\lesson{2}{Sep 15 2021 Wed (11:42:11)}{Poetry}{Unit 1}

When you want to analyze poetry, you need to consider the following elements:

\begin{itemize}
  \item Words
  \item Images
  \item Figurative Language
  \item Tone
  \item Structure
  \item Meaning
\end{itemize}

\subsection*{Different Parts of Speech}
\label{sub_sec:different_parts_of_speech}

\begin{definition}[Simile]
  Figure of speech that compares two unlike things using
  the words \textit{"like"} or \textit{"as"}
\end{definition}

\begin{definition}[Personification]
  Figure of speech in which inanimate or
  nonhuman things are given human characteristics or abilities
\end{definition}
\begin{definition}[Metaphor]
  Figure of speech that compares two unlike things without using any comparison
  words
\end{definition}

\begin{definition}[Apostrophe]
  Figure of speech that directly addresses an abstract quality, a nonhuman, or
  an individual that is not present
\end{definition}

\begin{definition}[Hyperbole]
  Figure of speech that uses exaggeration for effect
\end{definition}

\begin{definition}[Onomatopeia]
  Figure of speech in which the sound of the word corresponds to its meaning
\end{definition}

\begin{definition}[Alliteration]
  Repetition of initial consonant sounds
\end{definition}

\begin{definition}[Assonance]
  Repetition of vowel sounds within the word
\end{definition}

\begin{definition}[Consonance]
  Repetition of consonant sounds at the end of words
\end{definition}

\begin{definition}[Understatement]
  Figure of speech that makes something seem less important or serious to
  emphasize the opposite
\end{definition}

\begin{definition}[Rhyme]
  Identical sounds, usually at the end of words or lines of poetry where the
  final vowel sound and following consonant sounds are the same
\end{definition}

% subsection different_parts_of_speech (end)

\newpage
